% 原书第 212—214 页；第 13.6.2 小节全文。
\subsection{研究 2}
\label{subsec:ch13-study2}

在第二项模拟研究中，我们提高了数据生成过程的复杂性，并将样本量减小为 $n=30$。与前面一样，从均值为 $0$、方差为 $1$、相关结构相同的多元正态分布中，生成九个基线协变量。还生成一个二值变量 $R$，以概率 $0.5$ 等于 $1$，以概率 $0.5$ 等于 $-1$。最后一个协变量 $Z$，作为这些基线协变量和随机噪声 $U_Z$ 的函数生成：
\[
 Z=R\times\logit^{-1}(W1+W4+W7+0.5U_Z),
\]
其中，$U_Z$ 独立抽自标准正态分布。同样，干预 $A$ 采用平衡分配的随机化。对于配对匹配试验，使用非二分匹配算法 \texttt{nbpMatch}，考察两组匹配变量（Beck et al.，2016）。第一组按 $R$ 匹配，它是一个对 $Z$ 有较强影响的基线协变量；第二组按 $\{R,W2,W5,W8\}$ 匹配。随后，为每个单位生成结局
\[
 \begin{aligned}
 Y=\logit^{-1}\bigl[&0.75A+0.5(W2+W5+W8)+1.5Z+0.25U_Y\\
                  &{}+0.75A(W2-W5)+0.5AZ\bigr]/7.5,
 \end{aligned}
\]
其中，$U_Y$ 抽自标准正态分布。因此，结局为一个取值在 $[0,1]$ 内的连续变量，例如比例。还通过干预将 $A$ 设为 $a$，生成反事实结局 $Y_1$ 与 $Y_0$。对于每种研究设计，将这一数据生成过程重复 $2500$ 次。按照前面的方法计算 SATE 和 PATE。SATE 的范围为 $0.2\%$—$3.3\%$，均值为 $1.6\%$；PATE 的真值为 $1.6\%$。表 13.2 展示了基线协变量与结局之间的关系，以及自适应配对匹配方案。

\begin{table}[htbp]
 \centering
 \caption{模拟研究 2：协变量与结局的关系，以及自适应配对匹配方案}
 \label{tab:ch13-study2-relations}
 \begingroup
 \small
 \setlength{\tabcolsep}{4pt}
 \renewcommand{\arraystretch}{1.22}
 \begin{tabular}{lccccccccccc}
  \toprule
  &&\multicolumn{3}{c}{$\overbrace{\hspace{1.8cm}}^{\text{相关系数 }0.5}$}
   &\multicolumn{3}{c}{$\overbrace{\hspace{1.8cm}}^{\text{相关系数 }0.5}$}
   &\multicolumn{3}{c}{$\overbrace{\hspace{1.8cm}}^{\text{相关系数 }0}$}&\\
  &$R$&$W1$&$W2$&$W3$&$W4$&$W5$&$W6$&$W7$&$W8$&$W9$&$Z$\\
  \midrule
  协变量 $Z$ 的父节点&$\checkmark$&$\checkmark$&&&$\checkmark$&&&$\checkmark$&&&\\
  结局 $Y$ 的父节点&&&$\checkmark$&&&$\checkmark$&&&$\checkmark$&&$\checkmark$\\
  \midrule
  匹配集 1&$\checkmark$&&&&&&&&&&\\
  匹配集 2&$\checkmark$&&$\checkmark$&&&$\checkmark$&&&$\checkmark$&&\\
  \bottomrule
 \end{tabular}
 \endgroup
\end{table}

我们比较了相同的算法：未经调整的估计量、调整无关协变量 $W9$ 的 MLE、数据自适应估计条件结局均值的 TMLE，以及将条件结局均值的数据自适应估计与数据自适应靶向更新相结合的 C-TMLE。初始估计条件结局均值 $\bar Q_0(A,W)$ 的算法库，包含 $12$ 个 logistic 回归工作模型，每个模型都有截距、暴露 $A$ 的主效应，以及一个候选调整变量的主效应。候选调整变量为 $\{\varnothing,R,W1,\ldots,W9,Z\}$。协同估计暴露机制 $g_0(A\mid W)$ 的算法库，也包含 $12$ 个 logistic 回归工作模型，每个模型都有截距及一个候选调整变量的主效应，候选变量同样为 $\{\varnothing,R,W1,\ldots,W9,Z\}$。我们采用相同的表现度量和交叉验证方案。与前面一样，对于未经调整的估计量和 MLE，推断基于估计的影响曲线；对于 TMLE，则基于影响曲线的交叉验证估计，见第 13.5 节。我们假设标准化估计量服从 Student $t$ 分布：无匹配试验的自由度为 $n-2=28$，有匹配试验的自由度为 $n/2-1=14$。

\noindent\textbf{结果。}
第二项模拟研究的结果见表 13.3，总体上与上述发现一致。配对匹配即使只采用一个协变量，如匹配集 1，也提高了分析精度。以样本效应而非总体效应为目标，进一步提高了效率。在初始估计 $\bar Q_0(A,W)$ 时，纳入工作模型的数据自适应选择，获得了更高精度；最有效的分析采用 C-TMLE。事实上，无匹配试验中 PATE 的未经调整估计量，其 MSE 几乎是按预后协变量匹配，即匹配集 2 时，SATE 的 C-TMLE 的 $4.5$ 倍。这使检测干预效应的功效增加了 $29\%$。

\begin{table}[htbp]
 \centering
 \caption{模拟研究 2 中估计量表现的汇总}
 \label{tab:ch13-study2}
 \begingroup
 \small
 \setlength{\tabcolsep}{3pt}
 \renewcommand{\arraystretch}{1.16}
 \begin{tabular}{lrrrrrrrrrr}
  \toprule
  &\multicolumn{5}{c}{PATE}&\multicolumn{5}{c}{SATE}\\
  \cmidrule(lr){2-6}\cmidrule(lr){7-11}
  &MSE&rMSE&$\hat\sigma$&功效&覆盖率&MSE&rMSE&$\hat\sigma$&功效&覆盖率\\
  \midrule
  \multicolumn{11}{l}{\textbf{无匹配}}\\
  未调整&$1.8\mathrm E{-4}$&1.00&0.013&0.24&0.94&$1.6\mathrm E{-4}$&1.12&0.013&0.24&0.95\\
  MLE&$1.8\mathrm E{-4}$&0.95&0.012&0.25&0.93&$1.7\mathrm E{-4}$&1.06&0.012&0.25&0.94\\
  TMLE&$1.2\mathrm E{-4}$&1.50&0.010&0.33&0.94&$9.8\mathrm E{-5}$&1.79&0.010&0.33&0.96\\
  C-TMLE&$1.1\mathrm E{-4}$&1.54&0.010&0.34&0.93&$9.5\mathrm E{-5}$&1.85&0.010&0.34&0.96\\
  \addlinespace[.3em]
  \multicolumn{11}{l}{\textbf{匹配集 1}}\\
  未调整&$1.1\mathrm E{-4}$&1.54&0.012&0.21&0.98&$9.2\mathrm E{-5}$&1.90&0.011&0.28&0.97\\
  MLE&$1.2\mathrm E{-4}$&1.48&0.012&0.23&0.97&$9.7\mathrm E{-5}$&1.81&0.011&0.29&0.97\\
  TMLE&$9.2\mathrm E{-5}$&1.91&0.010&0.31&0.97&$6.9\mathrm E{-5}$&2.52&0.009&0.40&0.96\\
  C-TMLE&$9.0\mathrm E{-5}$&1.95&0.010&0.33&0.96&$6.9\mathrm E{-5}$&2.53&0.008&0.44&0.95\\
  \addlinespace[.3em]
  \multicolumn{11}{l}{\textbf{匹配集 2}}\\
  未调整&$6.5\mathrm E{-5}$&2.70&0.011&0.17&0.99&$4.6\mathrm E{-5}$&3.79&0.009&0.37&0.98\\
  MLE&$7.3\mathrm E{-5}$&2.41&0.011&0.20&0.99&$5.4\mathrm E{-5}$&3.27&0.009&0.37&0.98\\
  TMLE&$5.3\mathrm E{-5}$&3.30&0.009&0.28&0.99&$3.8\mathrm E{-5}$&4.66&0.008&0.47&0.98\\
  C-TMLE&$5.3\mathrm E{-5}$&3.28&0.009&0.32&0.99&$3.9\mathrm E{-5}$&4.44&0.007&0.53&0.97\\
  \bottomrule
 \end{tabular}
 \endgroup
 \par\smallskip
 \begin{minipage}{\linewidth}\footnotesize
 各行表示研究设计与估计量：未经调整的估计量、调整 $W9$ 的 MLE、数据自适应选择初始估计量的 TMLE，以及采用数据自适应选择的 C-TMLE。各列表示估计量表现：MSE 为偏差平方加方差；rMSE 为无匹配试验中 PATE 的未经调整估计量的 MSE，除以另一个估计量的 MSE；$\hat\sigma$ 为平均标准误估计；其余为功效和覆盖率。
 \end{minipage}
\end{table}

在这些模拟中，参数的指定对估计量表现产生了明显影响。先关注 PATE 的估计，再讨论 SATE。以总体效应为推断目标时，虽然配对匹配减小了 MSE，但它带来的实际功效改善却有所减弱。这很可能是因为，在无匹配试验中标准误被略微低估，而在配对匹配试验中则被高估。事实上，无匹配试验中，$95\%$ 置信区间的覆盖率略低于名义水平，为 $93\%$—$94\%$；而在匹配良好，即匹配集 2 时，覆盖率接近 $100\%$。在这组模拟中，用于配对匹配设计方差估计的修正因子 $\rho_n$，见式 (13.7)，约等于 $0$。因此，配对匹配试验中的方差估计量相当保守，交叉验证选择方案对无匹配试验的优化更充分。在无匹配研究中，有 $10\%$ 选择调整 $R$ 的 logistic 回归模型，来初始估计 $\bar Q_0(A,W)$；在对 $R$ 匹配良好，即匹配集 1 的研究中，这一比例为 $7\%$。此外，当按多个协变量匹配，即匹配集 2 时，$\bar Q_0(A,W)$ 的工作模型选择与无匹配试验非常相似。

相比之下，估计 SATE 时，更小的 MSE 转化为更高的实际功效，同时保持了名义水平、甚至偏于保守的置信区间覆盖率。例如，TMLE 在无匹配试验中的实际功效为 $33\%$，按单个协变量匹配时为 $40\%$，按多个协变量匹配时为 $47\%$。同样，C-TMLE 在无匹配试验中的实际功效为 $34\%$，按单个协变量配对匹配时为 $44\%$，按多个协变量匹配时为 $53\%$。在无匹配研究中，有 $10\%$ 选择调整 $R$ 的工作模型，来初始估计 $\bar Q_0(A,W)$；而在对 $R$ 匹配良好，即匹配集 1 的研究中，只有 $2\%$ 这样选择。在后一情形中，算法对其他预后基线协变量，如 $W2$ 和 $Z$，给予了更多权重。
\FloatBarrier

